\chapter{LITERATURE REVIEW}

\section{Introduction}
The evolution of dynamic data structures is a cornerstone of modern software 
engineering, fundamentally dictating how applications allocate, manage, and 
process memory at runtime. Dynamic arrays, owing to their contiguous memory 
layout and $O(1)$ random access, are among the most heavily utilized structures 
across computational domains. However, standardizing their resizing policies 
has historically involved rigid compromises between spatial efficiency (memory 
utilization) and temporal efficiency (computational overhead). 

This chapter critically examines the existing body of literature surrounding 
memory allocation, algorithmic complexity, optimal resizable arrays, and 
adaptive data structures. Section 2.2 reviews existing implementations of 
resizable arrays in programming languages and theoretical works from key 
researchers such as Brodnik, Sartor, and Bille. Section 2.3 identifies the 
research gaps, focusing on the lack of workload-aware and velocity-sensitive 
dynamic growth algorithms. Finally, Section 2.4 summarizes the key findings of 
the literature review.

\section{Existing Works}

\subsection{Standard Language Implementations (C++, Java, CPython)}
Dynamic arrays are standard components in the standard libraries of most 
mainstream programming languages. C++'s \texttt{std::vector}, Java's 
\texttt{ArrayList}, and Python's \texttt{list} all implement dynamic arrays, but 
they utilize different growth strategies that reflect varying design philosophies.

In the C++ Standard Template Library (STL), \texttt{std::vector} typically uses 
a growth factor of $g = 2.0$ in common implementations (such as GCC's libstdc++ 
and MSVC). When the vector is full, its capacity is doubled. This aggressive 
doubling ensures that the frequency of reallocation remains very low, which is 
highly desirable in performance-critical system programming. However, a growth 
factor of $2.0$ has severe memory implications. In the worst-case scenario, the 
vector can waste up to $50\%$ of its allocated memory. Mathematically, the 
space overhead is bounded by $O(N)$ where $N$ is the size of the array. Furthermore, 
a growth factor of $2.0$ prevents the reuse of previously deallocated memory 
blocks. Because the new capacity ($C_{\text{new}} = 2 C_{\text{old}}$) is always 
greater than the sum of all previous allocations ($C_0 + C_1 + ... + C_{k-1}$), 
the allocator cannot reuse the contiguous blocks that were recently freed, 
forcing the operating system to continuously locate new blocks higher up in the 
virtual memory address space.

Java's \texttt{ArrayList} utilizes a more conservative growth factor of $g = 1.5$. 
The new capacity is calculated as:
\begin{equation}
    C_{\text{new}} = C_{\text{old}} + (C_{\text{old}} \gg 1)
\end{equation}
By growing by $1.5\times$ instead of doubling, Java reduces the worst-case 
memory waste to $33.3\%$. Furthermore, a growth factor of $1.5$ theoretically 
allows for the reuse of previously freed memory blocks after a few reallocations, 
improving memory compaction. However, this savings in space comes at the price 
of time: a $1.5\times$ growth policy increases the number of reallocations 
and elements copied by approximately $1.7\times$ compared to a doubling policy, 
increasing garbage collection pressure and computational latency \cite{java_data_structures}.

CPython, the reference implementation of Python, implements its \texttt{list} 
structure using a highly conservative, size-dependent non-geometric growth 
factor. The allocation formula is designed to provide rapid growth for very small 
lists, and conservative growth for larger lists to conserve memory:
\begin{equation}
    C_{\text{new}} = C_{\text{old}} + (C_{\text{old}} \gg 3) + (C_{\text{old}} < 9 ? 3 : 6)
\end{equation}
For large arrays, the term $(C_{\text{old}} \gg 3)$ dominates, resulting in an 
asymptotic growth factor of $1.125$ ($1 \frac{1}{8}$). This highly conservative 
policy ensures that CPython lists waste very little memory (around $11.1\%$ 
asymptotically). However, the trade-off is a massive increase in reallocation 
frequency. For large scale data insertions, a growth factor of $1.125$ causes 
continuous memory reallocation and copying. In resource-constrained systems or 
high-performance computing, the CPU time spent in CPython's allocator during 
list appends represents a major bottleneck.

\subsection{Theoretical Work on Resizable Arrays}
To address the time-space trade-offs of simple geometric expansion, several 
researchers have proposed theoretical architectures for resizable arrays to analyze 
their asymptotic limits \cite{time_space_complexity}. 

Brodnik et al. \cite{optimal_resizable} proposed a landmark data structure 
representing an optimal resizable array. They proved that it is possible to 
maintain a resizable array that wastes only $O(\sqrt{N})$ space, while still 
maintaining $O(1)$ worst-case time complexity for lookups and insertions. Their 
design achieves this by dividing the array into a series of blocks of 
exponentially increasing sizes. Instead of maintaining a single contiguous block, 
the structure maintains a directory of pointers to blocks of size $2^{\lfloor i/2 \rfloor}$. 
While this theoretical design solves the geometric waste issue ($O(\sqrt{N})$ 
wasted space is asymptotically much smaller than $O(N)$), it introduces a level 
of indirection. Every lookup requires resolving a pointer through the directory, 
which bypasses the hardware prefetcher and causes L1/L2 cache misses, making the 
raw lookup speed significantly slower than a standard contiguous array.

Sartor \cite{memory_allocation} analyzed the memory allocation behaviors of 
dynamic arrays in virtual memory systems. Sartor highlighted that the virtual memory 
allocator's page size (typically 4KB) plays a crucial role in actual memory waste. 
If a dynamic array requests an allocation that does not align with page boundaries, 
the operating system allocates a whole page anyway, resulting in hidden 
fragmentation. Sartor proposed that array growth factors should align with 
underlying memory pages, suggesting that resizing algorithms should be aware of 
allocator-level metrics. However, Sartor's models remained static and did not 
adapt to runtime ingestion velocity.

Bille et al. \cite{resizable_optimal_time_space} extended the theoretical bounds 
of resizable arrays to parallel and concurrent environments. They proposed 
algorithms that achieve optimal space-time trade-offs under dynamic updates 
by amortizing the copying overhead across multiple threads. Similarly, research 
into extendible arrays \cite{allocating_storage, managing_storage} in database 
systems focused on managing secondary storage allocations. These works utilized 
index structures to coordinate multidimensional array growth, ensuring that 
data could be appended along multiple axes without copying the entire database. 
However, the mathematical and structural overhead of managing these index 
tables makes them impractical for standard, general-purpose programming 
languages where low latency and fast memory cache-locality are the highest 
priorities \cite{fast_array}.

\subsection{Specialized and Adaptive Memory Structures}
Another class of related work focuses on adaptive memory structures. The Adaptive 
Packed Memory Array (APMA) \cite{adaptive_packed} is a prime example. The APMA 
is designed to maintain sorted order among elements while allowing rapid 
insertions and deletions. It dynamically adjusts the density thresholds (packed 
fractions) in different segments of the array to control memory fragmentation 
and minimize element movements. While the APMA is highly adaptive, its design 
is optimized for maintaining sorted lists, which requires maintaining empty 
``gaps'' between elements. This is fundamentally different from a standard 
sequential dynamic array where elements are appended contiguously.

Additionally, researchers in signal processing have explored adaptive filtering and 
array processing algorithms \cite{trends_adaptive}. These systems adaptively 
tune filter coefficients based on input signal rates. While the domain is 
different, the concept of utilizing feedback loops (such as Kalman filters or 
recursive least squares) to adjust algorithmic parameters in real time serves 
as a conceptual inspiration for our workload-aware resizing algorithm.

Other specialized data structures, such as the Black-White Array \cite{black_white_array} 
and Z-rays \cite{z_rays}, focus on reducing the copying cost of dynamic arrays 
by using a dual-buffer system or virtual memory remapping. The Black-White 
Array splits the array into two segments to allow parallel insertions at both 
ends, while Z-rays utilize virtual memory page table manipulation to expand 
arrays without physically copying data. While remapping pages in virtual memory 
is highly efficient for extremely large arrays, it requires direct interaction 
with the operating system's kernel and is slow for small arrays due to system 
call overhead.

Furthermore, alternative designs include logarithmic architectures that scale 
the growth factor logarithmically to prevent large-scale overallocation at 
high capacities \cite{logarithmic_arrays}. For concurrent and persistent 
scenarios, lock-free vectors \cite{lock_free_vector} and fully persistent 
arrays \cite{persistent_arrays} allow threads to access historical versions 
of the array without copying data. For database query execution, 
block-structured arrays specialized for query processing 
\cite{query_processing} eliminate contiguous reallocations entirely.

\section{Research Gap}
Despite the extensive literature on resizable arrays, a significant research gap 
remains at the intersection of practical contiguous arrays and runtime 
adaptability. The existing work is sharply divided into two paradigms:
\begin{enumerate}
    \item \textbf{Rigid Contiguous Arrays:} The standard implementations (C++, Java, 
    Python) maintain a strictly contiguous block of memory to guarantee raw 
    $O(1)$ CPU-cache-efficient lookup speeds, but they rely on static growth 
    factors ($2.0\times$, $1.5\times$, $1.125\times$) that completely ignore the 
    nature of the workload.
    \item \textbf{Complex Non-Contiguous Structures:} Theoretical designs (e.g., 
    Brodnik's segmented arrays, packed memory arrays, virtual page remappers) 
    achieve better theoretical space bounds, but they destroy cache locality 
    due to pointer indirection, require kernel-level support, or add massive 
    CPU overhead for index calculations.
\end{enumerate}

Specifically, no existing contiguous dynamic array resizing algorithm incorporates:
\begin{itemize}
    \item \textbf{Workload-Aware Insertion Velocity ($v$):} Current systems 
    do not observe how fast elements are being appended. A fast, high-velocity 
    stream of data indicates an ongoing burst, which requires a large growth 
    factor to minimize frequent reallocations. A slow, trickle-like insertion 
    suggests that the process may be ending, requiring a conservative growth 
    factor to prevent leaving large amounts of memory unused.
    \item \textbf{Resize History Decoupling ($k$):} Although Python's resizing 
    has some size dependency, it does not decay the growth factor smoothly as 
    a function of the resize iteration $k$. As an array grows larger, the cost 
    of a resize grows linearly, making the penalty of overallocation increasingly 
    severe. None of the standard implementations decay the growth factor 
    harmonically to prevent massive absolute memory waste at scale.
    \item \textbf{Multi-Objective Parameter Optimization:} The parameters 
    of existing growth algorithms are chosen using heuristic rules-of-thumb 
    (e.g., doubling because it is simple). There has been no systematic attempt 
    to define a unified objective function $J$ that penalizes memory waste, 
    element copying, and runtime execution time, and to optimize this function 
    using modern hyper-parameter tuning techniques like Bayesian Optimization.
\end{itemize}

This research directly addresses this gap by formulating a contiguous dynamic 
array that maintains $O(1)$ random access, yet adaptively adjusts its growth 
factor at runtime using a harmonic decay of resize history $k$, scaled by the 
insertion velocity $v$.

\section{Summary for the Chapter}
This chapter reviewed the existing literature on resizable arrays, memory 
allocation, and adaptive data structures. Section 2.2 analyzed standard dynamic 
array implementations in mainstream programming languages (C++, Java, CPython), 
highlighting the rigid trade-offs between memory waste and allocation frequency. 
It also examined theoretical resizable array designs proposed by Brodnik, Sartor, 
and Bille, noting that while they achieve better space complexity, their 
pointer indirection degrades CPU cache performance. Section 2.3 identified a 
critical research gap: the lack of a contiguous, general-purpose dynamic 
array resizing algorithm that incorporates real-time insertion velocity and 
resize history feedback to dynamically tune the growth factor. This gap 
motivates the design of the workload-aware adaptive growth models detailed in 
the next chapter.
